\section{Introduction}
\label{sec:intro}

The onset of nucleate boiling (ONB) is the point at which the first vapour
bubbles nucleate and grow on a heated wall, marking the departure of a boiling
curve from its single-phase line. In forced-convection subcooled flow it
governs the lower bound of two-phase operation in microchannel electronics
cooling, the thermal-margin budget of nuclear fuel assemblies, and the design
of compact refrigerant evaporators. Predicting ONB accurately therefore has
direct engineering value: an over-conservative incipience estimate wastes
pumping power and surface area, while an under-conservative one risks premature
boiling, flow instability and, downstream, departure from nucleate boiling.

Classical ONB models predict the wall superheat $\dTonb = T_w - T_{\mathrm{sat}}$
at incipience from a force or thermal balance on a nucleating bubble. The Hsu
\cite{HSU1962} criterion and its descendants---Davis and Anderson
\cite{DAVIS1966}, Bergles and Rohsenow \cite{BERGLES1964}, Sato and Matsumura
\cite{SATOMATSUMURA1964}, and the mechanistic model of Basu et al. \cite{BASU2002}---each
embed assumptions (saturated bulk, a specific cavity-size spectrum, a fixed
thermal-layer thickness) calibrated within a narrow experimental window. As a
result they extrapolate poorly: applied outside their calibration range they
can err by an order of magnitude, and they disagree with one another by several
kelvin even on canonical water data. Modern applications, however, span a
parameter space far wider than any single correlation was built for---hydraulic
diameters from tens of micrometres to centimetres, pressures from atmospheric to
the nuclear range, mass fluxes over an order of magnitude, and both water and
low-GWP refrigerants---so a predictor that remains accurate and physically
consistent \emph{across} that space is needed.

Data-driven and physics-informed neural networks (PINNs) offer a route to such
breadth, but two obstacles recur in boiling. First, ONB data are scarce and
heterogeneous, scattered across decades of literature in incompatible units and
reporting conventions; a purely data-driven fit overfits and violates known
physical trends. Second, the relevant conservation laws (mass, momentum, energy
with phase change) are stiff and, near incipience, only weakly constrain the
scalar quantity of interest, so a full field-resolving PINN is both
data-hungry and ill-posed for sparse incipience data.

In a companion study (Phase~1) we showed that a \emph{physics-regularized}
network---one that predicts the scalar ONB targets directly while being
softly constrained by the Hsu nucleation criterion and physical
monotonicity---outperforms classical correlations for pool boiling and learns a
transferable representation of surface morphology (roughness, contact angle,
cavity spectrum). The present work extends that framework to subcooled forced
convection. The central question is whether the pool-boiling surface
representation can be \emph{reused} and combined with a flow representation to
predict ONB over the much larger forced-convection parameter space.

Our contributions are:
\begin{enumerate}
\item A surface- and flow-conditioned PINN that transfer-learns a frozen
pool-boiling surface encoder and adds a dedicated flow encoder, predicting both
$\dTonb$ and $\qonb$ across $\Dh = 0.35$--\SI{9.5}{mm} and
$P = 0.1$--\SI{16}{MPa}.
\item Identification and resolution of two systematic failure modes: a
subcooling--Hsu conflict at low superheat, fixed by imposing the Hsu criterion
as a one-sided \emph{lower bound}; and a pressure-insensitivity, fixed by adding
\emph{reduced pressure} as an explicit input feature so the network recovers the
$P\!\uparrow\!\Rightarrow\!\dTonb\!\downarrow$ trend.
\item A deep-ensemble uncertainty quantification with an honest calibration
analysis distinguishing epistemic from the dominant aleatoric scatter.
\item A curated, openly released dataset of 351 ONB points harmonised from nine
studies, with explicit removal of pool-boiling contamination.
\end{enumerate}

The remainder of the paper formulates the problem and architecture
(Section~\ref{sec:method}), describes the data (Section~\ref{sec:data}),
presents forward accuracy, physics consistency and uncertainty
(Section~\ref{sec:results}), discusses implications and limitations
(Section~\ref{sec:discussion}) and concludes (Section~\ref{sec:conclusions}).
